%% kapa: exam=1 round=35 year=2024 period=1 form=A subject=civil_law subject_ko=민법 seq=9 number=9
%% source: Q-Net 감정평가사 기출문제 (cst00309) artlSeq=5213229 file=raw/1008/5213229/1교시 문제지 원본.pdf
%% conversion: kordoc 4.24.51 -> tools/gen1.py (auto); review=proofread-sonnet; flags=none
\begin{question}{9}
불공정한 법률행위에 관한 설명으로 옳지 않은 것은? (다툼이 있으면 판례에 따름)
\begin{choices}
  \item 불공정한 법률행위에 해당하는지는 원칙적으로 법률행위 시를 기준으로 판단한다.
  \item 대리인에 의한 법률행위의 경우, 궁박 상태의 여부는 본인을 기준으로 판단한다.
  \item 경매에는 불공정한 법률행위에 관한 민법 제104조가 적용되지 않는다.
  \item 불공정한 법률행위는 추인으로 유효로 될 수 없지만 법정추인은 인정된다.
  \item 불공정한 법률행위는 이를 기초로 새로운 이해관계를 맺은 선의의 제3자에 대해서도 무효이다.
\end{choices}
\end{question}
\answer{4}
